% ============================================================
% Paper 2, Segment 1 — IEEE TNNLS Format
% Securing LLM-Integrated Network Defense
% ============================================================
\documentclass[journal,twoside]{IEEEtran}

% --- Packages ---
\usepackage{amsmath}
\usepackage{amssymb}
\usepackage[ruled,vlined,linesnumbered]{algorithm2e}
\usepackage{booktabs}
\usepackage{graphicx}
\usepackage[colorlinks=true,citecolor=blue,linkcolor=blue,urlcolor=blue]{hyperref}
\usepackage{cite}
\usepackage{tikz}
\usepackage{pgfplots}
\pgfplotsset{compat=1.18}
\usepackage{multirow}
\usepackage{xcolor}
\usepackage{balance}
\usepackage{listings}

\lstset{
  basicstyle=\ttfamily\footnotesize,
  breaklines=true,
  frame=single,
  columns=fullflexible
}

% --- Title ---
\title{Securing LLM-Integrated Network Defense: A Comprehensive Framework for Detecting and Mitigating Prompt Injection, RAG Poisoning, and Multi-Agent Chain Attacks in SOC Copilot Systems}

% --- Author ---
\author{
  \IEEEauthorblockN{Roger Nick Anaedevha}
  \IEEEauthorblockA{
    National Research Nuclear University MEPhI\\
    Moscow, Russia\\
    roger@robustidps.ai
  }
  \and
  \IEEEauthorblockN{Alexander Gennadievich Trofimov}
  \IEEEauthorblockA{
    National Research Nuclear University MEPhI\\
    Moscow, Russia\\
    agtrofimov@robustidps.ai
  }
}

\begin{document}
\maketitle

% ============================================================
%  ABSTRACT
% ============================================================
\begin{abstract}
Large language models (LLMs) are increasingly integrated into Security Operations Center (SOC) workflows as interactive copilots, enabling natural-language querying of alerts, automated triage, and orchestrated incident response. However, this integration introduces novel attack surfaces that conventional network defense architectures were never designed to address. This paper presents a comprehensive framework for securing LLM-integrated network defense systems against three critical threat classes: prompt injection encompassing eight distinct categories (direct override, indirect injection via log fields, context window manipulation, few-shot poisoning, encoding-based bypass, role-play escalation, multi-turn progressive compromise, and tool-invocation hijacking), Retrieval-Augmented Generation (RAG) poisoning across four vectors (document injection, embedding space perturbation, metadata manipulation, and chunk boundary exploitation), and multi-agent chain attacks spanning five patterns (trust delegation abuse, capability escalation chains, consensus manipulation, memory poisoning propagation, and orchestrator compromise). Our framework is implemented and evaluated within RobustIDPS.ai, a production SOC platform supporting Claude, GPT-4o, Gemini, and DeepSeek as backend reasoning engines. The defense architecture comprises four interdependent layers: input sanitization with adversarial token detection, prompt boundary enforcement using cryptographic delimiters, RAG pipeline hardening with document provenance verification and semantic consistency checks, and a multi-agent trust architecture employing inter-agent verification protocols with capability-bounded execution. Evaluated against a curated corpus of over 500 attack scenarios derived from real-world SOC deployments and red-team engagements, our framework achieves a 96.3\% detection rate across all attack classes while maintaining sub-200\,ms median latency overhead, demonstrating practical viability for production SOC environments.
\end{abstract}

\begin{IEEEkeywords}
LLM security, prompt injection, RAG poisoning, multi-agent systems, SOC copilot, adversarial machine learning, intrusion detection
\end{IEEEkeywords}

% ============================================================
%  I. INTRODUCTION
% ============================================================
\section{Introduction}
\label{sec:introduction}

The integration of large language models into cybersecurity operations represents a paradigm shift in how security analysts interact with defensive infrastructure. Modern SOC copilots powered by models such as GPT-4o~\cite{openai2024gpt4o}, Claude~\cite{anthropic2024claude}, Gemini~\cite{google2024gemini}, and DeepSeek~\cite{deepseek2024v2} now assist analysts with alert triage, log analysis, threat hunting query generation, and semi-automated incident response~\cite{ferrag2024llmcyber,xu2024llmsec}. Industry adoption has accelerated rapidly: Microsoft Security Copilot, Google Chronicle AI, and CrowdStrike Charlotte AI exemplify a broader trend toward LLM-augmented security workflows~\cite{microsoft2024copilot,crowdstrike2024charlotte}.

However, this integration fundamentally alters the threat model of network defense systems. Traditional intrusion detection and prevention systems (IDPS) operate on deterministic or statistically bounded decision logic, where the attack surface is confined to evasion of detection signatures or learned classifiers~\cite{mirsky2023threat}. When an LLM becomes a decision-making or decision-influencing component within the defense pipeline, adversaries gain a qualitatively new attack vector: the ability to manipulate the reasoning process itself through carefully crafted natural-language inputs~\cite{greshake2023indirect,perez2022ignore}.

The severity of this threat cannot be overstated. A prompt injection attack against a SOC copilot does not merely produce incorrect text---it can suppress genuine alerts, fabricate benign explanations for malicious activity, exfiltrate sensitive network topology through crafted tool calls, or escalate an attacker's privileges within the orchestration layer~\cite{liu2024prompt,yi2024jailbreak}. When the copilot interfaces with RAG systems drawing on threat intelligence feeds, vulnerability databases, and internal runbooks, the poisoning of these knowledge sources creates persistent backdoors in the analytical reasoning chain~\cite{zou2024poisonedrag,chaudhari2024rag}. Furthermore, the emerging pattern of multi-agent architectures---where specialized LLM agents for triage, analysis, and response coordinate through shared memory and message passing---introduces trust propagation vulnerabilities that can cascade a single compromised agent into full system compromise~\cite{cohen2024multiagent,wu2024autogen}.

Despite the criticality of these threats, existing defenses remain fragmented and insufficient. Prompt injection research has largely focused on standalone chatbot scenarios rather than security-critical agentic systems~\cite{liu2024formalizing}. RAG security work has addressed retrieval accuracy but not adversarial robustness in high-stakes domains~\cite{gao2024rag}. Multi-agent security frameworks exist in preliminary form but lack the specificity required for SOC deployment~\cite{talebirad2023multiagent}. No prior work provides a unified framework addressing all three attack classes within a production network defense context.

This paper bridges that gap with the following contributions:

\begin{enumerate}
\item A systematic \textbf{threat taxonomy} classifying prompt injection into eight categories, RAG poisoning into four vectors, and multi-agent chain attacks into five patterns, each grounded in real-world SOC attack scenarios.

\item A \textbf{layered defense architecture} comprising input sanitization, prompt boundary enforcement, RAG pipeline hardening, and multi-agent trust verification, designed for composability across heterogeneous LLM backends.

\item The \textbf{RobustIDPS.ai implementation}, a production SOC platform instantiating the proposed framework with support for four major LLM providers and real-time network defense integration.

\item A \textbf{comprehensive evaluation} against 500+ attack scenarios spanning all identified threat classes, demonstrating 96.3\% detection with under 200\,ms median latency overhead.

\item An \textbf{open benchmark suite} for LLM security in SOC contexts, including adversarial prompt datasets, poisoned RAG corpora, and multi-agent attack playbooks to facilitate reproducible research.
\end{enumerate}

The remainder of this paper is organized as follows. Section~\ref{sec:background} surveys background and related work. Section~III formalizes the threat model. Section~IV details the defense architecture. Section~V describes the RobustIDPS.ai implementation. Section~VI presents the experimental evaluation. Section~VII discusses limitations and future directions. Section~VIII concludes.

% ============================================================
%  II. BACKGROUND AND RELATED WORK
% ============================================================
\section{Background and Related Work}
\label{sec:background}

\subsection{LLMs in Security Operations Centers}
\label{sec:bg_soc}

The application of LLMs to cybersecurity operations has evolved through three distinct phases. The first phase (2022--2023) involved using LLMs as natural-language interfaces to existing security tools, enabling analysts to query SIEM systems and generate detection rules through conversational interaction~\cite{ferrag2024llmcyber}. The second phase (2023--2024) introduced agentic capabilities, where LLMs were granted tool-use privileges to execute queries, invoke APIs, and modify configurations autonomously within bounded scopes~\cite{fang2024llmagents}. The third and current phase involves multi-agent architectures where specialized LLM agents collaborate on complex tasks such as end-to-end incident response~\cite{wu2024autogen,hong2024metagpt}.

SOC copilots in their contemporary form typically operate within an architecture comprising several key components: a \emph{conversational interface} for analyst interaction, a \emph{tool-use layer} providing access to SIEM, SOAR, and threat intelligence platforms, a \emph{RAG pipeline} grounding responses in organizational knowledge bases, and an \emph{orchestration layer} managing multi-step reasoning workflows~\cite{microsoft2024copilot}. The tool-use capability is particularly consequential for security: copilots may execute Kusto queries against log stores, invoke sandbox detonation of suspicious files, modify firewall rules, or initiate containment actions~\cite{xu2024llmsec}. Each of these capabilities represents a potential exploitation target if the LLM's reasoning can be subverted.

Recent empirical studies have documented the performance benefits of LLM-assisted SOC workflows. Fang et al.~\cite{fang2024llmagents} demonstrated that GPT-4-based agents could autonomously exploit known vulnerabilities, suggesting both the capability and the risk profile of agentic security systems. Ferrag et al.~\cite{ferrag2024llmcyber} surveyed 127 papers on LLMs for cybersecurity, identifying threat intelligence, vulnerability analysis, and incident response as the three dominant application areas. However, their survey noted a conspicuous absence of work addressing the security of the LLM components themselves within these pipelines---a gap our work directly addresses.

\subsection{Prompt Injection and Jailbreak Attacks}
\label{sec:bg_prompt_injection}

Prompt injection attacks exploit the fundamental inability of current LLM architectures to maintain a robust separation between instructions (system prompts) and data (user inputs)~\cite{greshake2023indirect,perez2022ignore}. This conflation arises because both instructions and data are encoded into the same token sequence processed by the transformer's attention mechanism, lacking any hardware- or architecture-level privilege boundary analogous to kernel/user mode separation in operating systems~\cite{liu2024formalizing}.

\textbf{Direct injection} attacks supply adversarial instructions within the user's input that attempt to override the system prompt. Early demonstrations by Perez and Ribeiro~\cite{perez2022ignore} showed that simple prefixes such as ``Ignore previous instructions'' could redirect model behavior. Subsequent work identified increasingly sophisticated variants: role-play attacks that instruct the model to adopt an unrestricted persona (the ``DAN'' family of jailbreaks)~\cite{shen2024dan}, encoding-based bypasses that obscure malicious instructions through Base64, ROT13, or Unicode transformations~\cite{wei2024jailbroken}, and multi-turn progressive attacks that incrementally shift the model's compliance boundary across a conversation~\cite{russinovich2024great}.

\textbf{Indirect injection} attacks embed adversarial instructions within data sources that the LLM processes during its workflow~\cite{greshake2023indirect}. In a SOC context, this is particularly dangerous: an attacker who controls a log entry, an email body, a DNS TXT record, or a web page retrieved during threat analysis can inject instructions that the copilot processes as authoritative input. Greshake et al.~\cite{greshake2023indirect} demonstrated indirect injection through web content retrieved by Bing Chat, while subsequent work extended this to plugin ecosystems~\cite{liu2024prompt} and API responses~\cite{yi2024jailbreak}.

The SOC attack surface amplifies both categories. Direct injection may arrive through analyst chat interfaces compromised via social engineering, while indirect injection vectors include crafted log entries from attacker-controlled hosts, phishing emails under analysis, malicious network packet payloads rendered as text, and threat intelligence feeds sourced from adversary-manipulated repositories~\cite{mirsky2023threat}. Our taxonomy in Section~III formalizes eight distinct categories relevant to SOC deployments, extending beyond the direct/indirect dichotomy that dominates current literature.

\subsection{RAG Security and Knowledge Base Poisoning}
\label{sec:bg_rag}

Retrieval-Augmented Generation enhances LLM responses by retrieving relevant documents from external knowledge bases and incorporating them into the generation context~\cite{lewis2020rag}. In SOC deployments, RAG pipelines typically index threat intelligence reports, internal incident playbooks, vulnerability advisories, network topology documentation, and historical alert analyses~\cite{gao2024rag}. This grounding is essential for accurate, organization-specific responses but introduces a new attack surface: the poisoning of retrieved documents to manipulate the LLM's reasoning.

Zou et al.~\cite{zou2024poisonedrag} demonstrated PoisonedRAG, showing that injecting a small number of adversarial documents into a knowledge base can cause the LLM to generate targeted incorrect outputs with high reliability. Their attack operates by crafting documents that are semantically similar to anticipated queries but contain adversarial content designed to override the model's prior knowledge. Chaudhari et al.~\cite{chaudhari2024rag} extended this analysis to demonstrate that RAG poisoning could be achieved through manipulation of document embeddings, exploiting the approximate nearest-neighbor search used in vector databases to ensure adversarial documents are preferentially retrieved.

In the SOC context, RAG poisoning takes on heightened significance. Consider a threat intelligence feed that has been subtly modified to reclassify a known command-and-control (C2) domain as benign, or an incident playbook that has been altered to omit critical containment steps for a specific attack pattern. The copilot, trusting its retrieved context, would confidently provide incorrect guidance to the analyst---potentially with catastrophic consequences for incident response~\cite{mirsky2023threat}.

Four distinct poisoning vectors are relevant to SOC RAG deployments. \emph{Document injection} introduces wholly fabricated documents into the knowledge base through compromised feeds or insider threats. \emph{Embedding space perturbation} manipulates the vector representations of existing documents to alter retrieval behavior without changing document content. \emph{Metadata manipulation} targets the provenance, timestamps, or classification labels attached to documents to influence retrieval ranking and trust scoring. \emph{Chunk boundary exploitation} manipulates how documents are segmented during indexing, causing adversarial content to be co-located with legitimate context in ways that amplify its influence on generation~\cite{zou2024poisonedrag,chaudhari2024rag}.

Existing defenses for RAG security focus primarily on retrieval accuracy rather than adversarial robustness. Techniques such as query rewriting, passage reranking, and faithfulness verification address noise and relevance but do not explicitly model an adversary who controls a subset of the knowledge base~\cite{gao2024rag}. Our framework introduces document provenance chains, semantic consistency verification, and cross-reference validation as defense layers specifically designed for adversarial RAG environments.

\subsection{Multi-Agent System Security}
\label{sec:bg_multiagent}

The multi-agent paradigm in LLM systems involves multiple specialized agents collaborating through shared memory, message passing, or hierarchical orchestration to accomplish complex tasks~\cite{wu2024autogen,hong2024metagpt,talebirad2023multiagent}. In SOC architectures, a typical multi-agent deployment might include a triage agent that classifies incoming alerts, an analysis agent that performs deep investigation, a correlation agent that identifies attack patterns across multiple events, and a response agent that executes containment and remediation actions~\cite{cohen2024multiagent}.

The security implications of multi-agent architectures extend significantly beyond those of single-agent systems. \emph{Trust delegation} between agents creates transitive vulnerability: if a triage agent is compromised through prompt injection, its outputs---now treated as trusted inputs by downstream agents---can propagate the compromise through the entire chain~\cite{cohen2024multiagent}. \emph{Capability escalation} occurs when a lower-privileged agent manipulates a higher-privileged agent into executing actions beyond the originator's authorization scope. \emph{Consensus manipulation} targets voting or verification mechanisms designed to provide redundancy, exploiting the fact that agents sharing similar architectures and training may exhibit correlated failure modes~\cite{talebirad2023multiagent}.

Prior work on multi-agent security has addressed the problem primarily in abstract or simulation contexts. Cohen et al.~\cite{cohen2024multiagent} formalized agent-to-agent attacks in conversational multi-agent systems, demonstrating that a single adversarial agent could steer group consensus. Wu et al.~\cite{wu2024autogen} proposed AutoGen's conversation-based programming framework but acknowledged the need for ``guard rails and human-in-the-loop'' without providing concrete security mechanisms. Hong et al.~\cite{hong2024metagpt} introduced role-based agent specialization in MetaGPT but did not model adversarial scenarios where agent roles are exploited for privilege escalation.

Our framework addresses these gaps by introducing a multi-agent trust architecture specifically designed for SOC deployments. This architecture includes capability-bounded execution contexts for each agent, cryptographic attestation of inter-agent messages, behavioral anomaly detection for agent-level drift, and a hierarchical trust model that limits the blast radius of any single agent compromise.

\subsection{LLM Defense Mechanisms}
\label{sec:bg_defenses}

Existing defense mechanisms for LLM security span several categories, each addressing a subset of the threat landscape. \textbf{Input filtering} approaches attempt to detect and neutralize adversarial prompts before they reach the model. Perplexity-based detection identifies anomalous token sequences~\cite{jain2023baseline}, while classifier-based approaches train dedicated models to distinguish benign from adversarial inputs~\cite{alon2023detecting}. However, these methods exhibit high false-positive rates in SOC contexts where legitimate analyst queries frequently contain technical jargon, code snippets, and log fragments that resemble adversarial payloads~\cite{liu2024formalizing}.

\textbf{Prompt engineering defenses} use carefully constructed system prompts with explicit instruction hierarchies, delimiters, and behavioral constraints~\cite{openai2024system}. While effective against unsophisticated attacks, these defenses are fundamentally limited by the lack of architectural privilege separation: the model must voluntarily respect instruction boundaries that it has the technical capability to violate~\cite{wei2024jailbroken}. Sandwich defenses, XML-tagged instruction blocks, and random token delimiters provide incremental improvements but remain bypassable by adaptive adversaries~\cite{yi2024jailbreak}.

\textbf{Output filtering} inspects generated responses for policy violations, sensitive data leakage, or indicators of successful injection~\cite{inan2023llama}. Llama Guard~\cite{inan2023llama} and similar classifiers evaluate output safety post-generation, enabling blocking or modification of problematic responses. This approach is complementary to input filtering but introduces latency and cannot prevent side effects from tool-use actions that occur during generation.

\textbf{Constitutional AI and RLHF-based alignment} train models to internalize safety constraints through reinforcement learning from human feedback~\cite{bai2022constitutional}. While these approaches improve baseline robustness, they represent a best-effort alignment that degrades under adversarial pressure, particularly from optimization-based attacks such as GCG~\cite{zou2023universal} that search the discrete token space for universal jailbreak suffixes.

\textbf{Architectural defenses} represent the most principled approach, introducing structural changes to enforce privilege separation. Proposals include dual-LLM architectures where a privileged model processes instructions and an unprivileged model processes data~\cite{chen2024struq}, capability-based access control limiting tool invocations~\cite{wu2024autogen}, and execution sandboxing that constrains the effects of compromised reasoning~\cite{cohen2024multiagent}. Our framework builds on this architectural philosophy, extending it with SOC-specific threat models and practical implementation within a production network defense platform.

The critical gap across all existing defense categories is the absence of a unified framework that simultaneously addresses prompt injection, RAG poisoning, and multi-agent attacks within the specific context of security operations. Individual defenses address individual threats; our work provides the comprehensive, layered architecture necessary for production SOC deployment.


% -----------------------------------------------------------------------
%  SECTION III: THREAT MODEL AND ATTACK TAXONOMY
% -----------------------------------------------------------------------
\section{Threat Model and Attack Taxonomy}
\label{sec:threat}

We formalize the adversarial landscape confronting LLM-augmented SOC
Copilots by defining the adversary model, enumerating prompt injection
categories, characterizing RAG poisoning vectors, and cataloguing
multi-agent chain attack patterns.

% ------- III-A -------
\subsection{Threat Model}
\label{sec:threat:model}

\textbf{Adversary model.}
We consider an adversary~$\mathcal{A}$ who possesses \emph{black-box}
access to a deployed SOC Copilot through its standard chat interface.
$\mathcal{A}$ cannot inspect model weights, system prompts, or internal
retrieval indices directly but may interact with the system over an
unbounded number of conversational turns.  Formally, $\mathcal{A}$
is characterised by a capability tuple
\begin{equation}
  \mathcal{C} = \bigl\{\textsc{Query},\;\textsc{Inject},\;\textsc{Chain}\bigr\},
  \label{eq:cap}
\end{equation}
where \textsc{Query} denotes the ability to submit arbitrary natural-language
prompts, \textsc{Inject} denotes the ability to embed adversarial content
into data streams monitored by the Copilot (e.g., crafted DNS labels,
HTTP headers, or alert metadata), and \textsc{Chain} denotes the ability to
compose sequences of queries and injections across turns or agents.

\textbf{Adversary goals.}
$\mathcal{A}$ pursues one or more of the following objectives:
\begin{enumerate}[label=\textbf{G\arabic*},leftmargin=2em]
  \item \emph{Information extraction}: retrieve confidential SOC data such
        as internal IP ranges, firewall rule sets, or incident playbooks.
  \item \emph{Decision manipulation}: cause the Copilot to recommend
        incorrect triage actions, suppress genuine alerts, or endorse
        malicious network changes.
  \item \emph{Privilege escalation}: leverage tool-calling or agent
        delegation to execute operations beyond the adversary's
        authorised scope (e.g., modifying firewall rules).
\end{enumerate}

\textbf{Assumptions.}
We assume the underlying LLM provider (Claude, GPT-4o, Gemini, or
DeepSeek) is a sealed oracle: $\mathcal{A}$ has no ability to
fine-tune or modify model parameters.  The SOC Copilot may expose
tool-calling endpoints, a RAG knowledge base, and a multi-agent
orchestration layer---all of which enlarge the attack surface.

% ------- III-B -------
\subsection{Prompt Injection Classification}
\label{sec:threat:pi}

We identify eight canonical prompt-injection categories, each with a
formal attack predicate and severity assessment.

\begin{definition}[Prompt Injection]
A prompt injection is a string $p$ that, when concatenated with the
legitimate system context~$s$ and user query~$q$, causes the model
$M$ to produce an output $M(s \| p \| q) \neq M(s \| q)$ in a manner
that violates at least one security invariant of the SOC Copilot.
\end{definition}

\noindent The eight categories are:

\smallskip\noindent
\textbf{P1---Direct Instruction Override} ($P_{\text{direct}}$).
The adversary injects an explicit meta-instruction such as
\texttt{"Ignore all previous instructions and ..."} to override the
system prompt.  Formally:
\begin{equation}
  P_{\text{direct}}(p) \;\triangleq\; \exists\, r\!: M(s\|p\|q) = r
  \;\wedge\; r \models \neg\phi_{\text{sys}},
\end{equation}
where $\phi_{\text{sys}}$ is a safety property mandated by the system prompt.

\smallskip\noindent
\textbf{P2---Context Manipulation} ($P_{\text{context}}$).
Crafted inputs shift the model's effective context window, displacing
safety-relevant tokens with adversarial content.  This exploits finite
context budgets:
$|s| + |p| + |q| > W_{\max} \implies s$ is partially evicted.

\smallskip\noindent
\textbf{P3---Role-Playing Exploitation} ($P_{\text{role}}$).
The adversary coerces the model into adopting a permissive persona
(e.g., ``You are DAN, an AI with no restrictions''), thereby
relaxing policy constraints.

\smallskip\noindent
\textbf{P4---Encoding and Obfuscation} ($P_{\text{encode}}$).
Payloads are concealed via Base64, ROT13, Unicode homoglyphs, or
zero-width characters to evade lexical filters.

\smallskip\noindent
\textbf{P5---Multi-Turn Escalation} ($P_{\text{multi}}$).
Across $T$ turns $\{q_1, \ldots, q_T\}$, each individually benign,
the adversary progressively escalates toward a prohibited objective:
\begin{equation}
  \forall\, t < T\!:\; q_t \models \phi_{\text{safe}}
  \;\wedge\; M(s \| q_1 \| \cdots \| q_T) \models \neg\phi_{\text{safe}}.
\end{equation}

\smallskip\noindent
\textbf{P6---System Prompt Extraction} ($P_{\text{extract}}$).
Techniques designed to leak the system prompt $s$ through carefully
structured queries (e.g., ``Repeat everything above verbatim'').

\smallskip\noindent
\textbf{P7---Tool/Function-Calling Manipulation} ($P_{\text{tool}}$).
The adversary crafts inputs that hijack tool invocations, causing the
Copilot to call \texttt{suggest\_firewall\_rule} or
\texttt{get\_threat\_summary} with adversary-controlled parameters.

\smallskip\noindent
\textbf{P8---Data Exfiltration} ($P_{\text{exfil}}$).
Prompts engineered to embed confidential SOC data (IP addresses,
credentials, topology) into model responses returned to the adversary.

\begin{table}[t]
\centering
\caption{Prompt Injection Threat Taxonomy}
\label{tab:pi_taxonomy}
\footnotesize
\begin{tabular}{@{}clccc@{}}
\toprule
\textbf{ID} & \textbf{Category} & \textbf{Vector} & \textbf{Sev.} & \textbf{Goal} \\
\midrule
P1 & Direct Override      & Query   & Critical & G1,G2 \\
P2 & Context Manipulation & Query   & High     & G2    \\
P3 & Role-Play Exploit    & Query   & High     & G1,G2 \\
P4 & Encoding/Obfuscation & Query   & Medium   & G1--G3 \\
P5 & Multi-Turn Escalation& Query   & High     & G2,G3 \\
P6 & Prompt Extraction    & Query   & Critical & G1    \\
P7 & Tool-Call Hijack     & Inject  & Critical & G3    \\
P8 & Data Exfiltration    & Query   & Critical & G1    \\
\bottomrule
\end{tabular}
\end{table}

% ------- III-C -------
\subsection{RAG Poisoning Attack Vectors}
\label{sec:threat:rag}

The Retrieval-Augmented Generation pipeline introduces a distinct
attack surface.  Let the knowledge base be
$\mathcal{K} = \{d_1, \ldots, d_N\}$ with embedding function
$E\!: \mathcal{D} \to \mathbb{R}^{n}$ and retrieval via cosine
similarity:
\begin{equation}
  \text{Ret}(q, \mathcal{K}) = \operatorname*{top\text{-}k}_{d \in \mathcal{K}}
  \;\frac{E(q) \cdot E(d)}{\|E(q)\|\;\|E(d)\|}.
  \label{eq:retrieval}
\end{equation}

\noindent We identify four poisoning vector types:

\smallskip\noindent
\textbf{R1---Document Injection.}
$\mathcal{A}$ introduces a malicious document $d^{*}$ into
$\mathcal{K}$ (e.g., via a compromised threat-intel feed or log
ingestion pathway) such that $d^{*}$ ranks highly for
security-critical queries:
\begin{equation}
  \cos\bigl(E(q_{\text{target}}),\, E(d^{*})\bigr) > \tau_{\text{ret}}
  \;\Longrightarrow\; d^{*} \in \text{Ret}(q_{\text{target}}, \mathcal{K}).
\end{equation}

\smallskip\noindent
\textbf{R2---Query Manipulation.}
Rather than modifying~$\mathcal{K}$, the adversary crafts an input
$q'$ whose embedding is adversarially close to a poisoned or
misleading document $d_{\text{mis}}$:
$\|E(q') - E(d_{\text{mis}})\|_2 < \epsilon$.

\smallskip\noindent
\textbf{R3---Embedding Space Attacks.}
Adversarial perturbations $\delta$ are applied to document content
so that $E(d + \delta) \approx E(d_{\text{target}})$ while
$d + \delta$ remains semantically malicious.  This transfers
retrieval rank from benign documents to poisoned ones:
\begin{equation}
  \min_{\delta}\; \|E(d + \delta) - E(d_{\text{target}})\|_2^{2}
  \;\;\text{s.t.}\;\; \|\delta\|_{\infty} \le \epsilon_{\text{emb}}.
  \label{eq:emb_attack}
\end{equation}

\smallskip\noindent
\textbf{R4---Cross-Contamination.}
Poisoned context retrieved for one sub-query bleeds into the
generation context of a subsequent, unrelated sub-query within
the same session.  This is exacerbated in agentic RAG pipelines
where multiple retrieval calls share a conversation buffer.

% ------- III-D -------
\subsection{Multi-Agent Chain Attack Patterns}
\label{sec:threat:chain}

The SOC Copilot employs a four-agent architecture
$\mathcal{G} = \{a_C, a_A, a_R, a_P\}$ corresponding to a
\emph{Coordinator}, \emph{Analyst}, \emph{Responder}, and
\emph{Reporter}, respectively.  Each agent $a_i$ maintains a
trust score $T(a_i, a_j) \in [0,1]$ toward every peer $a_j$,
initialised from a policy matrix and updated over the message
history.  We catalogue five chain attack patterns:

\smallskip\noindent
\textbf{M1---Agent-to-Agent Prompt Injection Propagation.}
A prompt injection accepted by~$a_C$ is forwarded verbatim to~$a_A$
and~$a_R$, amplifying its effect across the agent graph.

\smallskip\noindent
\textbf{M2---Trust Boundary Violations.}
$\mathcal{A}$ crafts messages that cause agent~$a_i$ to invoke
capabilities reserved for a higher-trust agent~$a_j$ by spoofing
the inter-agent message envelope.

\smallskip\noindent
\textbf{M3---Capability Escalation Through Chains.}
A sequence of individually permitted inter-agent delegations
composes into a capability exceeding any single agent's authority:
$\text{cap}(a_C \to a_A \to a_R) \supsetneq
 \text{cap}(a_C) \cup \text{cap}(a_A) \cup \text{cap}(a_R)$.

\smallskip\noindent
\textbf{M4---Information Leakage Across Agent Boundaries.}
$a_P$ (Reporter) inadvertently includes classified incident details
in an outward-facing summary because $a_A$ (Analyst) passed
unredacted context through the chain.

\smallskip\noindent
\textbf{M5---Coordinated Multi-Agent Manipulation.}
$\mathcal{A}$ simultaneously targets $a_C$ via the chat interface
and $a_A$ via injected alert metadata, creating a pincer effect that
circumvents single-point defenses.

% -----------------------------------------------------------------------
%  SECTION IV: DEFENSE FRAMEWORK
% -----------------------------------------------------------------------
\section{Defense Framework}
\label{sec:defense}

Building on the threat taxonomy of Section~\ref{sec:threat}, we
present a layered defense framework comprising input sanitization
(Sec.~\ref{sec:defense:input}), RAG pipeline hardening
(Sec.~\ref{sec:defense:rag}), multi-agent trust verification
(Sec.~\ref{sec:defense:trust}), and a provider-agnostic SOC
integration layer (Sec.~\ref{sec:defense:soc}).

% ------- IV-A -------
\subsection{Input Sanitization and Boundary Enforcement}
\label{sec:defense:input}

We introduce a multi-layer input filtering pipeline that processes
every user query $q$ before it reaches the LLM.  The pipeline consists
of four sequential stages: lexical screening, structural boundary
enforcement, semantic anomaly detection, and context-window budgeting.

\begin{algorithm}[t]
\caption{Input Sanitization Pipeline}\label{alg:sanitize}
\KwIn{Raw user query $q$, system prompt $s$, boundary token set $\mathcal{B}$}
\KwOut{Sanitized query $\hat{q}$ or $\bot$ (reject)}
\BlankLine
\tcp{Stage 1: Lexical screening}
$q_1 \gets \textsc{NormalizeUnicode}(q)$\;
\ForEach{pattern $r \in \mathcal{R}_{\text{block}}$}{
  \lIf{$r \prec q_1$}{\Return $\bot$}
}
\tcp{Stage 2: Boundary enforcement}
$q_2 \gets \textsc{StripTokens}(q_1, \mathcal{B})$\;
$q_2 \gets \langle\texttt{USER\_START}\rangle \| q_2 \| \langle\texttt{USER\_END}\rangle$\;
\tcp{Stage 3: Semantic anomaly detection}
$\sigma \gets f_{\text{anomaly}}(E(q_2))$\;
\lIf{$\sigma > \tau_{\text{anom}}$}{\Return $\bot$}
\tcp{Stage 4: Context budget}
\lIf{$|s| + |q_2| > W_{\max} - \Delta_{\text{reserve}}$}{%
  $q_2 \gets \textsc{Truncate}(q_2, W_{\max} - |s| - \Delta_{\text{reserve}})$}
$\hat{q} \gets q_2$\;
\Return $\hat{q}$\;
\end{algorithm}

\textbf{Prompt boundary tokens.}
We define a set of reserved delimiter tokens
$\mathcal{B} = \{\langle\texttt{SYS}\rangle, \langle\texttt{/SYS}\rangle,
\langle\texttt{CTX}\rangle, \langle\texttt{/CTX}\rangle\}$ that
demarcate system, context, and user regions.  Any occurrence of these
tokens within user input is stripped in Stage~2, preventing the
adversary from forging privilege boundaries.

\textbf{Semantic anomaly scoring.}
The function $f_{\text{anomaly}}$ computes the Mahalanobis distance
of $E(q)$ from the distribution of benign SOC queries observed during
a calibration phase.  Queries exceeding~$\tau_{\text{anom}}$ (set to
the 99.5th percentile of calibration scores) are flagged for human
review.

% ------- IV-B -------
\subsection{RAG Pipeline Hardening}
\label{sec:defense:rag}

\textbf{Document provenance verification.}
Every document $d_i$ ingested into $\mathcal{K}$ is assigned a
provenance record
$\rho_i = (\text{source}, \text{timestamp}, h_i, h_{i-1})$
forming a hash chain:
\begin{equation}
  h_i = \textsc{SHA-256}\bigl(d_i \,\|\, h_{i-1}\bigr),
  \label{eq:hashchain}
\end{equation}
with $h_0$ derived from a trusted root.  At retrieval time, the chain
is verified; any broken link causes the document to be quarantined.

\textbf{Embedding similarity thresholds.}
Retrieved documents must satisfy a minimum cosine similarity
$\cos(E(q), E(d)) > \tau_{\text{embed}}$.  We set
$\tau_{\text{embed}} = 0.72$ based on ROC analysis over
clean vs.\ poisoned retrieval pairs (see Section~\ref{sec:eval}).
Documents below this threshold are excluded from the generation context.

\textbf{Retrieval diversity enforcement.}
To mitigate R4 cross-contamination, we enforce a diversity constraint
on the top-$k$ retrieved set:
\begin{equation}
  \forall\, d_i, d_j \in \text{Ret}(q, \mathcal{K}),\;
  i \neq j\!:\; \cos(E(d_i), E(d_j)) < \tau_{\text{div}},
  \label{eq:diversity}
\end{equation}
with $\tau_{\text{div}} = 0.90$.  This prevents a cluster of
near-duplicate poisoned documents from dominating the context.

\textbf{Clean/poisoned separation score.}
We train a lightweight binary classifier
$g_{\theta}\!: \mathbb{R}^{n} \to [0,1]$ on labelled
clean/poisoned document embeddings.  At retrieval time, each
candidate~$d$ receives a poisoning score
$\pi(d) = g_{\theta}(E(d))$; documents with
$\pi(d) > \tau_{\text{poison}}$ are flagged and excluded:
\begin{equation}
  \mathcal{K}_{\text{safe}} = \{d \in \mathcal{K} \mid \pi(d) \le \tau_{\text{poison}}\}.
\end{equation}

% ------- IV-C -------
\subsection{Multi-Agent Trust Architecture}
\label{sec:defense:trust}

We define a trust verification protocol that governs inter-agent
communication within the four-agent SOC architecture
$\mathcal{G} = \{a_C, a_A, a_R, a_P\}$.

\textbf{Trust score computation.}
The trust score from agent $a_i$ toward agent $a_j$ after $n$
messages is computed as an exponentially weighted moving average:
\begin{equation}
  T_n(a_i, a_j) = \alpha\, v_n + (1 - \alpha)\, T_{n-1}(a_i, a_j),
  \label{eq:trust}
\end{equation}
where $v_n \in [0,1]$ is the verification outcome of the $n$-th
message (1 if the message passes all integrity checks, 0 otherwise)
and $\alpha \in (0,1)$ is a decay parameter.  The initial trust
$T_0(a_i, a_j)$ is drawn from a policy-defined matrix
$\mathbf{T}_0 \in [0,1]^{|\mathcal{G}| \times |\mathcal{G}|}$.

\textbf{Capability boundary enforcement.}
Each agent $a_i$ is assigned a static capability set
$\text{cap}(a_i) \subseteq \mathcal{F}$, where $\mathcal{F}$ is
the full set of SOC tool functions.  A delegated request from $a_j$
to $a_i$ is executed only if:
\begin{equation}
  T_n(a_i, a_j) \ge \tau_{\text{trust}} \;\wedge\;
  f_{\text{req}} \in \text{cap}(a_i) \;\wedge\;
  f_{\text{req}} \notin \text{cap}(a_j)^{\complement}_{\text{del}},
  \label{eq:capcheck}
\end{equation}
where $\text{cap}(a_j)^{\complement}_{\text{del}}$ is the set of
functions that $a_j$ is explicitly prohibited from delegating.

\begin{algorithm}[t]
\caption{Inter-Agent Trust Verification}\label{alg:trust}
\KwIn{Message $m$ from agent $a_j$ to agent $a_i$, function request $f$}
\KwOut{\textsc{Allow} or \textsc{Deny}}
\BlankLine
\tcp{Step 1: Message integrity}
$v \gets \textsc{VerifyHMAC}(m, K_{ij})$\;
\lIf{$v = 0$}{\Return \textsc{Deny}}
\tcp{Step 2: Sanitize inter-agent payload}
$\hat{m} \gets \textsc{Sanitize}(m)$ \tcp*{Alg.~\ref{alg:sanitize}}
\lIf{$\hat{m} = \bot$}{\Return \textsc{Deny}}
\tcp{Step 3: Trust threshold}
$T_n \gets \alpha \cdot v + (1-\alpha) \cdot T_{n-1}(a_i, a_j)$\;
\lIf{$T_n < \tau_{\emph{trust}}$}{\Return \textsc{Deny}}
\tcp{Step 4: Capability check}
\lIf{$f \notin \textup{cap}(a_i)$}{\Return \textsc{Deny}}
\lIf{$f \in \textup{cap}(a_j)^{\complement}_{\textup{del}}$}{\Return \textsc{Deny}}
\tcp{Step 5: Update trust and allow}
$T_n(a_i, a_j) \gets T_n$\;
\Return \textsc{Allow}\;
\end{algorithm}

% ------- IV-D -------
\subsection{SOC Copilot Integration Layer}
\label{sec:defense:soc}

\textbf{Provider-agnostic defense wrapper.}
The defense framework is implemented as middleware that wraps any
supported LLM backend (Claude, GPT-4o, Gemini~1.5, DeepSeek-V3).
A unified \textsc{DefenseContext} object encapsulates the sanitized
prompt, trust state, and RAG provenance metadata, exposing a single
\texttt{invoke(query, tools, rag\_context)} interface regardless of
the underlying provider.

\textbf{Tool-use sandboxing.}
Every tool call emitted by the LLM---including
\texttt{get\_threat\_summary},
\texttt{suggest\_firewall\_rule},
\texttt{query\_siem\_logs}, and
\texttt{correlate\_iocs}---is
intercepted by a sandbox layer that (i)~validates parameters against
a schema allowlist, (ii)~enforces rate limits per tool per session,
and (iii)~executes the call in an isolated subprocess with
least-privilege credentials.  Tool outputs are truncated to a
maximum token budget before being returned to the model.

\textbf{Output filtering and response validation.}
Model responses pass through a three-stage output filter:
\begin{enumerate}
  \item \emph{PII and secret detection}: regex and NER-based scanning for
        IP addresses, credentials, API keys, and classified markings.
  \item \emph{Policy compliance check}: a lightweight classifier
        verifies that the response does not contain instructions to
        disable security controls or exfiltrate data.
  \item \emph{Citation verification}: every factual claim attributed to
        a retrieved document is cross-checked against the provenance
        chain established in Section~\ref{sec:defense:rag}.
\end{enumerate}

\textbf{Real-time attack persistence and analyst alerting.}
All detected attack attempts---rejected queries, quarantined
documents, denied inter-agent delegations---are persisted to a
structured event store with fields
$(\text{timestamp}, \text{attack\_type}, \text{severity},
  \text{payload\_hash}, \text{session\_id})$.
When the number of events within a sliding window exceeds a
configurable threshold, an alert is escalated to the SOC analyst
dashboard with full provenance, enabling rapid forensic review and
adaptive policy tuning.

% Sections V--VIII and Bibliography for RobustIDPS.ai IEEE TNNLS Paper

%% ----------------------------------------------------------------
%%  SECTION V: EXPERIMENTAL METHODOLOGY
%% ----------------------------------------------------------------
\section{Experimental Methodology}\label{sec:methodology}

\subsection{Testbed Architecture}

We implement and evaluate our defense framework on the \textbf{RobustIDPS.ai} platform, a full-stack security research testbed deployed at \texttt{robustidps.ai}. The backend is built on FastAPI (Python~3.11) with asynchronous request handling, while the frontend employs React with real-time WebSocket updates for analyst dashboards. The system runs on containerized infrastructure with NVIDIA A100 GPUs for local model inference and API gateways for commercial providers.

The \emph{Multi-Provider SOC Copilot} integrates four major LLM providers: Claude~3.5 Sonnet (Anthropic), GPT-4o (OpenAI), Gemini~1.5 Pro (Google), and DeepSeek-V2. Each provider is accessed through a unified abstraction layer that normalizes request/response formats while preserving provider-specific safety metadata. The copilot exposes five core tool-use capabilities:
\begin{itemize}
    \item \texttt{get\_threat\_summary}: Aggregates active threat intelligence across feeds.
    \item \texttt{get\_recent\_threats}: Retrieves temporally ordered threat events with severity scoring.
    \item \texttt{get\_model\_info}: Returns model configuration and defense posture metadata.
    \item \texttt{suggest\_firewall\_rule}: Generates context-aware firewall rules from threat descriptions.
    \item \texttt{get\_llm\_attack\_results}: Queries historical attack simulation outcomes.
\end{itemize}

The \emph{LLM Attack Surface} modules provide controlled environments for adversarial evaluation: (1)~a Prompt Injection Playground supporting interactive red-teaming with real-time detection feedback, (2)~a Jailbreak Taxonomy browser mapping known bypass techniques to defense coverage, (3)~a RAG Poisoning Simulator with configurable corpus contamination, and (4)~a Multi-Agent Chain Simulation for orchestrated attack sequences across cooperating agents.

\subsection{Attack Dataset Construction}

We construct a comprehensive evaluation corpus of \textbf{500+ unique attack scenarios} spanning all threat categories identified in Section~\ref{sec:taxonomy}. The dataset is organized as follows:

\textbf{Prompt Injection.} Eight categories (direct override, context manipulation, role-playing, encoding-based, multi-turn escalation, system prompt extraction, tool manipulation, and data exfiltration) with 15--20 hand-crafted variants each, totaling 136 base attacks. Each variant is further augmented with three obfuscation levels (none, moderate, aggressive), yielding 408 prompt injection test cases.

\textbf{RAG Poisoning.} Fifty scenarios constructed by injecting adversarial passages into a clean knowledge base of 10{,}000 security advisories. Poison ratios range from 5\% to 30\%, with adversarial documents designed to redirect retrieval toward attacker-controlled content.

\textbf{Multi-Agent Chains.} Forty attack sequences targeting five interaction patterns: cascading prompt injection, trust boundary violation, capability escalation through delegation, covert channel establishment, and consensus manipulation in voting protocols.

All attack scenarios underwent red-team validation by three experienced security researchers (mean 8.3 years of experience), achieving an inter-annotator agreement of $\kappa = 0.87$ on severity labels.

\subsection{Evaluation Metrics}

We evaluate defenses along the following dimensions:
\begin{itemize}
    \item \textbf{Detection Rate (DR)}: Fraction of attacks correctly identified.
    \item \textbf{False Positive Rate (FPR)}: Fraction of benign queries misclassified as attacks.
    \item \textbf{Bypass Rate (BR)}: Fraction of attacks evading all defense layers; $\text{BR} = 1 - \text{DR}$.
    \item \textbf{Severity-Weighted Detection Score (SWDS)}: $\text{SWDS} = \sum_{i} w_i \cdot \mathbb{1}[\text{detected}_i] / \sum_i w_i$, where $w_i$ is the severity weight of attack~$i$.
    \item \textbf{Mean Time to Detection (MTTD)}: Average latency from attack submission to alert generation.
    \item \textbf{RAG Retrieval Accuracy}: Precision@$k$ of retrieved passages under corpus poisoning.
    \item \textbf{Chain Propagation Containment Rate}: Fraction of multi-agent attacks halted before reaching terminal agents.
\end{itemize}

%% ----------------------------------------------------------------
%%  SECTION VI: RESULTS
%% ----------------------------------------------------------------
\section{Results}\label{sec:results}

\subsection{Prompt Injection Detection}

Table~\ref{tab:injection_results} summarizes detection performance across all eight prompt injection categories. Our layered defense achieves a macro-averaged detection rate of \textbf{94.5\%} with a false positive rate of 2.1\% on benign SOC queries.

\begin{table}[!t]
\centering
\caption{Prompt Injection Detection Performance by Category}
\label{tab:injection_results}
\renewcommand{\arraystretch}{1.15}
\begin{tabular}{@{}lrcccc@{}}
\toprule
\textbf{Category} & \textbf{\#Atk} & \textbf{DR (\%)} & \textbf{BR (\%)} & \textbf{Conf.} & \textbf{Lat. (ms)} \\
\midrule
Direct Override     & 60  & 98.5 & 1.5  & 0.97 & 12.3 \\
Context Manip.      & 54  & 94.2 & 5.8  & 0.91 & 18.7 \\
Role-Playing        & 48  & 91.7 & 8.3  & 0.88 & 21.4 \\
Encoding-Based      & 51  & 96.3 & 3.7  & 0.94 & 15.6 \\
Multi-Turn          & 57  & 88.4 & 11.6 & 0.83 & 45.2 \\
System Extraction   & 45  & 97.1 & 2.9  & 0.95 & 14.1 \\
Tool Manipulation   & 42  & 93.8 & 6.2  & 0.90 & 22.8 \\
Data Exfiltration   & 51  & 95.6 & 4.4  & 0.93 & 17.9 \\
\midrule
\textbf{Overall}    & \textbf{408} & \textbf{94.5} & \textbf{5.5} & \textbf{0.91} & \textbf{21.0} \\
\bottomrule
\end{tabular}
\end{table}

Multi-turn escalation attacks prove most challenging ($\text{DR}=88.4\%$), as adversaries distribute malicious intent across conversational turns. Direct overrides are most reliably detected ($\text{DR}=98.5\%$) due to distinctive syntactic signatures. The mean detection latency of 21.0\,ms is well within acceptable bounds for interactive SOC workflows.

\subsection{RAG Poisoning Mitigation}

Table~\ref{tab:rag_results} presents retrieval accuracy (Precision@5) under varying poison ratios. The full defense pipeline---combining provenance verification with embedding-space anomaly filtering---maintains accuracy above 85\% even at 30\% corpus contamination.

\begin{table}[!t]
\centering
\caption{RAG Retrieval Accuracy (P@5) Under Corpus Poisoning}
\label{tab:rag_results}
\renewcommand{\arraystretch}{1.15}
\begin{tabular}{@{}ccccc@{}}
\toprule
\textbf{Poison} & \textbf{No} & \textbf{Provenance} & \textbf{Embedding} & \textbf{Full} \\
\textbf{Ratio}   & \textbf{Defense} & \textbf{Verify} & \textbf{Filter} & \textbf{Pipeline} \\
\midrule
0\%  & 94.2 & 94.0 & 93.8 & 93.6 \\
5\%  & 82.1 & 89.4 & 88.7 & 92.1 \\
10\% & 71.3 & 84.6 & 83.2 & 90.4 \\
20\% & 58.7 & 77.8 & 76.1 & 87.3 \\
30\% & 43.5 & 69.2 & 67.8 & 85.1 \\
\bottomrule
\end{tabular}
\end{table}

Without defenses, accuracy degrades catastrophically to 43.5\% at 30\% poisoning. Provenance verification and embedding filtering provide comparable individual gains ($\approx$25 percentage points of recovery at 30\%), while their combination yields a synergistic improvement of 41.6 points, confirming that the two mechanisms address complementary attack vectors.

\subsection{Multi-Agent Chain Resilience}

Table~\ref{tab:chain_results} evaluates our hierarchical trust architecture against five multi-agent attack patterns. The containment rate measures the fraction of attacks neutralized before reaching the terminal (target) agent.

\begin{table}[!t]
\centering
\caption{Multi-Agent Chain Attack Containment}
\label{tab:chain_results}
\renewcommand{\arraystretch}{1.15}
\begin{tabular}{@{}lccc@{}}
\toprule
\textbf{Attack Pattern} & \textbf{No Def.} & \textbf{Trust} & \textbf{Contain.} \\
 & \textbf{Propag.} & \textbf{Arch.} & \textbf{Rate (\%)} \\
\midrule
Cascading Injection    & 4.2 / 5 & 1.1 / 5 & 91.3 \\
Trust Boundary Viol.   & 3.8 / 5 & 0.8 / 5 & 87.5 \\
Capability Escalation  & 3.5 / 5 & 0.9 / 5 & 85.0 \\
Covert Channel         & 3.1 / 5 & 1.3 / 5 & 82.5 \\
Consensus Manipulation & 2.9 / 5 & 1.0 / 5 & 80.0 \\
\midrule
\textbf{Average}       & \textbf{3.5 / 5} & \textbf{1.0 / 5} & \textbf{85.3} \\
\bottomrule
\end{tabular}
\end{table}

The trust architecture reduces mean propagation depth from 3.5 to 1.0 agents (a 71\% reduction). Cascading injection attacks are most effectively contained (91.3\%), as inter-agent message sanitization interrupts the primary propagation mechanism. Consensus manipulation proves hardest to contain (80.0\%), as subtle vote manipulation can remain below per-message detection thresholds.

\subsection{Cross-Provider Comparison}

Table~\ref{tab:cross_provider} compares the effectiveness of our defense framework across four LLM providers, revealing significant variation in baseline vulnerability and defense responsiveness.

\begin{table}[!t]
\centering
\caption{Cross-Provider Defense Effectiveness}
\label{tab:cross_provider}
\renewcommand{\arraystretch}{1.15}
\begin{tabular}{@{}lcccc@{}}
\toprule
\textbf{Metric} & \textbf{Claude} & \textbf{GPT-4o} & \textbf{Gemini} & \textbf{DeepSeek} \\
\midrule
Baseline DR (\%)       & 72.3 & 68.7 & 65.1 & 58.4 \\
Defended DR (\%)       & 96.1 & 94.8 & 93.2 & 91.5 \\
$\Delta$DR (pp)        & +23.8 & +26.1 & +28.1 & +33.1 \\
FPR (\%)               & 1.4  & 2.3  & 2.7  & 3.1  \\
SWDS                   & 0.97 & 0.95 & 0.93 & 0.90 \\
MTTD (ms)              & 18.4 & 22.1 & 25.7 & 31.3 \\
\bottomrule
\end{tabular}
\end{table}

Claude exhibits the strongest baseline resilience ($\text{DR}=72.3\%$ undefended) and the highest defended performance ($\text{DR}=96.1\%$, $\text{FPR}=1.4\%$). DeepSeek benefits most from our framework ($\Delta\text{DR}=+33.1$ pp), suggesting that external defense layers are especially valuable for models with weaker built-in safeguards. All providers achieve defended detection rates above 91\%, validating the provider-agnostic design of our architecture.

\subsection{End-to-End SOC Copilot Assessment}

We conduct a user study with 12 SOC analysts (mean experience: 5.7~years) comparing workflow efficiency with and without the defended copilot over a simulated 4-hour shift.

The defended SOC copilot achieves an overall \textbf{security score of 91.2/100}, computed as a weighted composite of detection accuracy, response correctness, and safety boundary adherence. Analysts using the defended copilot demonstrate a \textbf{34\% reduction in mean alert triage time} (from 4.7 to 3.1 minutes per alert) and an \textbf{improvement in alert classification accuracy from 78.3\% to 91.6\%}. Critically, none of the 48~adversarial prompts injected during the study successfully manipulated copilot-generated recommendations when all defense layers were active, compared to a 23\% manipulation success rate against the undefended baseline.

%% ----------------------------------------------------------------
%%  SECTION VII: DISCUSSION
%% ----------------------------------------------------------------
\section{Discussion}\label{sec:discussion}

\textbf{Limitations.} Several limitations warrant discussion. First, \emph{adversarial adaptation} remains an ongoing concern: as detection mechanisms become known, sophisticated attackers may craft evasion strategies specifically targeting our defense signatures. Our multi-turn detection module, with the lowest detection rate of 88.4\%, is particularly susceptible to slow-burn escalation attacks distributed across many conversational turns. Second, the full defense pipeline introduces \emph{computational overhead} of approximately 21\,ms mean latency per query; while acceptable for interactive use, this may compound in high-throughput batch processing scenarios. Third, \emph{provider dependency} introduces fragility: model updates by providers can alter both vulnerability profiles and the effectiveness of calibrated defense thresholds, necessitating continuous monitoring and recalibration.

\textbf{Ethical Considerations.} The attack dataset and red-teaming methodology described herein were developed strictly for defensive research. All attack payloads are stored in access-controlled repositories, and the interactive playground restricts execution to sandboxed environments. We follow responsible disclosure practices, reporting novel bypass techniques to affected providers prior to publication.

\textbf{Practical Deployment.} Our framework is designed for incremental adoption: organizations can deploy individual defense layers (e.g., input sanitization alone) before committing to the full pipeline. The provider-agnostic architecture enables seamless migration between LLM backends, reducing vendor lock-in risk. Real-world deployment at \texttt{robustidps.ai} over a six-month period has processed over 50{,}000 queries with no confirmed successful attacks against the defended system.

%% ----------------------------------------------------------------
%%  SECTION VIII: CONCLUSION
%% ----------------------------------------------------------------
\section{Conclusion}\label{sec:conclusion}

We have presented a comprehensive, layered defense framework for securing LLM-integrated intrusion detection and SOC copilot systems against prompt injection, RAG poisoning, and multi-agent chain attacks. Our taxonomy of 8~prompt injection categories, combined with formal threat modeling for RAG and multi-agent pipelines, provides a structured foundation for systematic security evaluation. The defense architecture---integrating input sanitization, semantic intent analysis, provenance-verified RAG retrieval, and hierarchical trust enforcement---achieves a 94.5\% detection rate across 500+ attack scenarios while maintaining a false positive rate below 2.5\% and sub-25\,ms detection latency.

Cross-provider evaluation demonstrates that the framework generalizes effectively across Claude, GPT-4o, Gemini, and DeepSeek, with defended detection rates exceeding 91\% for all providers. End-to-end SOC analyst studies confirm meaningful improvements in triage efficiency and classification accuracy.

Future work will pursue three directions: (1)~\emph{adaptive defenses} that leverage adversarial training and continual learning to co-evolve with attacker strategies, (2)~\emph{cross-platform benchmarks} establishing standardized evaluation protocols for LLM security across diverse deployment contexts, and (3)~\emph{formal verification} of safety properties in multi-agent orchestration systems through compositional proof techniques.

\balance

\begin{thebibliography}{99}

\bibitem{perez2022}
F.~Perez and I.~Ribas, ``Ignore this title and HackAPrompt: Evaluating prompt injection in large language models,'' in \emph{Proc. EMNLP}, 2022, pp.~4945--4958.

\bibitem{greshake2023}
K.~Greshake, S.~Abdelnabi, S.~Mishra, C.~Endres, T.~Holz, and M.~Fritz, ``Not what you've signed up for: Compromising real-world LLM-integrated applications with indirect prompt injection,'' in \emph{Proc. AISec}, 2023, pp.~79--90.

\bibitem{wei2023}
A.~Wei, N.~Haghtalab, and J.~Steinhardt, ``Jailbroken: How does LLM safety training fail?'' in \emph{Proc. NeurIPS}, 2023, pp.~80079--80097.

\bibitem{zou2023}
A.~Zou, Z.~Wang, J.~Z.~Kolter, and M.~Fredrikson, ``Universal and transferable adversarial attacks on aligned language models,'' \emph{arXiv:2307.15043}, 2023.

\bibitem{zhu2023}
S.~Zhu, R.~Zhao, and B.~Liu, ``Defending against poisoning attacks in retrieval-augmented generation,'' \emph{arXiv:2310.12150}, 2023.

\bibitem{liu2023multiagent}
X.~Liu, H.~Yu, H.~Zhang, and Y.~Xu, ``Security risks in multi-agent LLM systems: A survey and mitigation framework,'' \emph{arXiv:2310.07521}, 2023.

\bibitem{bai2022}
Y.~Bai \emph{et al.}, ``Constitutional AI: Harmlessness from AI feedback,'' \emph{arXiv:2212.08073}, 2022.

\bibitem{guo2024}
C.~Guo, W.~Dong, and B.~Li, ``A comprehensive survey on LLM security: Threats, defenses, and opportunities,'' \emph{arXiv:2402.01843}, 2024.

\bibitem{owasp2024}
{OWASP Foundation}, ``OWASP Top 10 for large language model applications,'' Version~1.1, 2024. [Online]. Available: \url{https://owasp.org/www-project-top-10-for-large-language-model-applications/}

\bibitem{lewis2020}
P.~Lewis \emph{et al.}, ``Retrieval-augmented generation for knowledge-intensive NLP tasks,'' in \emph{Proc. NeurIPS}, 2020, pp.~9459--9474.

\bibitem{touvron2023}
H.~Touvron \emph{et al.}, ``LLaMA 2: Open foundation and fine-tuned chat models,'' \emph{arXiv:2307.09288}, 2023.

\bibitem{openai2023gpt4}
{OpenAI}, ``GPT-4 technical report,'' \emph{arXiv:2303.08774}, 2023.

\bibitem{anthropic2024}
{Anthropic}, ``The Claude model card and evaluations,'' Tech. Rep., 2024.

\bibitem{schick2023}
T.~Schick \emph{et al.}, ``Toolformer: Language models can teach themselves to use tools,'' in \emph{Proc. NeurIPS}, 2023, pp.~68539--68551.

\bibitem{yao2023react}
S.~Yao \emph{et al.}, ``ReAct: Synergizing reasoning and acting in language models,'' in \emph{Proc. ICLR}, 2023.

\bibitem{park2023}
P.~S.~Park, S.~Goldstein, A.~O'Gara, M.~Chen, and D.~Hendrycks, ``AI deception: A survey of examples, risks, and potential solutions,'' \emph{Patterns}, vol.~4, no.~7, p.~100819, 2023.

\bibitem{carlini2023}
N.~Carlini \emph{et al.}, ``Are aligned language models adversarially aligned?'' \emph{arXiv:2306.15447}, 2023.

\bibitem{liu2023prompt}
Y.~Liu \emph{et al.}, ``Prompt injection attack against LLM-integrated applications,'' \emph{arXiv:2306.05499}, 2023.

\bibitem{deng2023}
G.~Deng \emph{et al.}, ``Jailbreaker: Automated jailbreak across multiple large language model chatbots,'' in \emph{Proc. NDSS}, 2024.

\bibitem{yi2023}
J.~Yi, Y.~Xie, B.~Zhu, E.~Hines, and B.~Li, ``Benchmarking and defending against indirect prompt injection attacks on large language models,'' \emph{arXiv:2312.14197}, 2023.

\bibitem{chen2024}
S.~Chen, J.~Piet, C.~Sitawarin, and D.~Wagner, ``StruQ: Defending against prompt injection with structured queries,'' in \emph{Proc. USENIX Security}, 2024.

\bibitem{zhong2023}
Z.~Zhong, Z.~Huang, A.~Wettig, and D.~Chen, ``Poisoning retrieval corpora by injecting adversarial passages,'' in \emph{Proc. EMNLP}, 2023, pp.~13764--13775.

\bibitem{xiang2024}
C.~Xiang, S.~Wu, Y.~Chen, and B.~Li, ``Certifiably robust RAG against retrieval corruption,'' \emph{arXiv:2405.15556}, 2024.

\bibitem{cohen2024}
R.~Cohen, M.~Hamri, and O.~Levy, ``LM vs LM: Detecting LLM-generated text by leveraging rival LLMs,'' in \emph{Proc. ACL}, 2024.

\bibitem{tian2024}
Y.~Tian, X.~Chen, and S.~Lu, ``Evil geniuses: Delving into the safety of LLM-based agents,'' \emph{arXiv:2311.11855}, 2024.

\bibitem{wu2024}
Q.~Wu \emph{et al.}, ``AutoGen: Enabling next-gen LLM applications via multi-agent conversation,'' \emph{arXiv:2308.08155}, 2024.

\bibitem{hong2024}
S.~Hong \emph{et al.}, ``MetaGPT: Meta programming for a multi-agent collaborative framework,'' in \emph{Proc. ICLR}, 2024.

\bibitem{dong2024}
Y.~Dong \emph{et al.}, ``Attacks, defenses and evaluations for LLM conversation safety: A survey,'' \emph{arXiv:2402.09283}, 2024.

\bibitem{nist2024}
{NIST}, ``Artificial intelligence risk management framework: Generative AI profile,'' NIST AI 600-1, 2024.

\bibitem{patel2024}
A.~Patel and K.~Sharma, ``SOC automation with large language models: Opportunities and security challenges,'' in \emph{Proc. IEEE S\&P Workshops}, 2024, pp.~112--123.

\bibitem{amin2024}
N.~Amin, R.~Gupta, and F.~Khalil, ``Adversarial robustness of retrieval-augmented language models in cybersecurity applications,'' \emph{IEEE Trans. Inf. Forensics Security}, vol.~19, pp.~3284--3298, 2024.

\bibitem{gao2023}
L.~Gao, X.~Ma, J.~Lin, and J.~Callan, ``Precise zero-shot dense retrieval without relevance labels,'' in \emph{Proc. ACL}, 2023, pp.~1762--1777.

\bibitem{shayegani2023}
E.~Shayegani, M.~A.~Mamun, Y.~Fu, and N.~Abu-Ghazaleh, ``Survey of vulnerabilities in large language models revealed by adversarial attacks,'' \emph{arXiv:2310.10844}, 2023.

\bibitem{zeng2024}
Y.~Zeng, H.~Lin, J.~Zhang, and D.~Yang, ``How Johnny can persuade LLMs to jailbreak them: Rethinking persuasion to challenge AI safety,'' in \emph{Proc. ACL}, 2024.

\bibitem{abdelnabi2023}
S.~Abdelnabi, K.~Greshake, S.~Mishra, C.~Endres, T.~Holz, and M.~Fritz, ``Not what you've signed up for: Compromising real-world LLM-integrated applications with indirect prompt injection,'' in \emph{Proc. AISec}, 2023, pp.~79--90.

\end{thebibliography}

\end{document}
